\chapter{HIV Transmission and Prevention}

\section*{Definition and Overview}
HIV transmission refers to the ways the human immunodeficiency virus passes from one person to another, and prevention refers to the highly effective strategies that stop it. HIV is transmitted through certain body fluids: blood, semen, vaginal fluids, rectal fluids, and breast milk. It is not spread by casual contact, air, or insect bites. Prevention combines condoms, pre-exposure prophylaxis (PrEP), post-exposure prophylaxis (PEP), safe injecting practices, testing, and treatment. The modern insight that an undetectable viral load means the virus cannot be transmitted sexually has transformed prevention.

\section*{Epidemiology}
New HIV diagnoses in Australia have fallen by around a third over the past decade, and transmission among men who have sex with men has declined sharply with PrEP and treatment. Mother-to-child transmission has been almost eliminated in Australia, occurring in fewer than one in a hundred births to women with HIV. Around 30,000 Australians live with HIV, and over 90\% have an undetectable viral load. Globally, transmission continues in regions with limited access to prevention, though new infections have fallen by around 40\% since the peak.

\section*{Aetiology and Pathophysiology}
\begin{itemize}
    \item \textbf{Routes of transmission:} unprotected vaginal and anal sex, sharing needles, mother-to-child, and contaminated blood or needles
    \item \textbf{Body fluids:} the virus is present in blood, semen, vaginal and rectal fluids, and breast milk
    \item \textbf{Risk factors:} the highest risks are receptive anal sex without condoms, sharing injecting equipment, and occupational needlesticks
    \item \textbf{Mother-to-child:} transmission can occur in pregnancy, during birth, or through breastfeeding, but is almost eliminated with treatment
    \item \textbf{Not transmitted by:} kissing, hugging, sharing food, toilets, swimming pools, coughing, or mosquito bites
    \item \textbf{Undetectable equals untransmittable:} effective treatment makes sexual transmission impossible
\end{itemize}

\section*{Clinical Features}
\begin{itemize}
    \item Transmission often occurs from people unaware they have HIV, which is why testing matters
    \item The viral load is highest in early infection, when transmission risk is greatest
    \item Untreated HIV has high transmission risk; treated HIV with undetectable viral load has zero sexual transmission risk
    \item Condoms are the most accessible barrier against HIV and other infections
    \item PrEP prevents acquisition in people at ongoing risk
    \item PEP prevents infection when started within 72 hours of exposure
    \item Safe injecting practices prevent transmission among people who inject drugs
\end{itemize}

\section*{Differential Diagnosis}
\begin{itemize}
    \item Risk assessment distinguishes ongoing risk, potential recent exposure, and known infection, which have different prevention strategies
    \item Other sexually transmitted infections share routes and are screened together
    \item People with undetectable viral load need no additional transmission precautions with partners
    \item Occupational exposures are managed through specific protocols
    \item PrEP, PEP, and treatment are selected according to the individual's situation
\end{itemize}

\section*{When to Seek Medical Care}
Anyone at risk should be tested regularly, and anyone with a potential exposure within the past 72 hours needs PEP urgently. PrEP is available from GPs and sexual health clinics.

\textbf{Call 000 (or local emergency services) for:}
\begin{itemize}
    \item An occupational needlestick or significant exposure requiring urgent assessment
    \item A sexual assault, which requires urgent PEP and other care
    \item Any severe acute illness within weeks of a possible exposure
    \item Symptoms suggesting acute HIV: high fever, rash, and severe sore throat with recent risk
\end{itemize}

\section*{Diagnosis and Investigations}
\begin{itemize}
    \item \textbf{HIV testing:} available through GPs, sexual health clinics, and home tests; most tests detect infection within 4 weeks
    \item \textbf{Baseline assessment for PrEP:} HIV test, kidney function, and hepatitis screening
    \item \textbf{Assessment for PEP:} confirmation of exposure and timing within 72 hours
    \item \textbf{STI screening:} syphilis, gonorrhoea, chlamydia, and hepatitis are checked together
    \item \textbf{Pregnancy testing:} where relevant
    \item \textbf{Follow-up testing:} after PEP, at 6 weeks and 3 months
\end{itemize}

\section*{Management}
\begin{itemize}
    \item \textbf{Condoms:} highly effective when used consistently and correctly
    \item \textbf{PrEP:} daily or event-based tenofovir-based tablets reduce acquisition risk by over 99\% when taken as prescribed
    \item \textbf{PEP:} 28 days of antiretroviral tablets started within 72 hours of exposure, with near-complete prevention
    \item \textbf{Testing and treatment:} diagnosing and treating everyone with HIV prevents onward transmission
    \item \textbf{Undetectable equals untransmittable:} counselling that effective treatment eliminates sexual transmission
    \item \textbf{Safe injecting:} never share equipment; use needle and syringe programs
    \item \textbf{Mother-to-child prevention:} antenatal testing and treatment
    \item \textbf{Partner notification:} supports testing and prevention
\end{itemize}

\section*{Complications}
\begin{itemize}
    \item HIV acquisition and the lifelong need for treatment if prevention fails
    \item Transmission from undiagnosed infection during the highly infectious early phase
    \item Co-infection with other sexually transmitted infections
    \item Occupational exposure in health workers
    \item Stigma and discrimination around testing and status
    \item Side effects and adherence requirements of PrEP and PEP
    \item Gaps in access to prevention for some communities
\end{itemize}

\section*{Prognosis and Outcomes}
HIV is almost entirely preventable. Consistent condom use, PrEP for those at risk, safe injecting, and universal treatment make acquisition and transmission rare. Australia has committed to ending HIV transmission, and current trends support this goal: with PrEP, prompt PEP, and treatment as prevention, the tools exist to eliminate new infections. The outlook for individuals using prevention is excellent, and for those who do acquire HIV, early treatment restores a normal life expectancy.

\section*{Prevention and Self-Care}
\begin{itemize}
    \item Use condoms for vaginal and anal sex
    \item Ask about PrEP if you are at ongoing risk
    \item Act fast after a possible exposure: PEP must start within 72 hours
    \item Get tested regularly and encourage partners to test
    \item Never share needles, syringes, or injecting equipment
    \item People with HIV take treatment to reach an undetectable viral load
    \item Pregnant women attend antenatal testing
    \item Talk openly with health professionals about risk
\end{itemize}

\section*{Sources and Further Reading}
The following reputable sources provide further information for healthcare students and patients seeking reliable, current advice about HIV transmission and prevention.

\begin{itemize}
    \item ASHM. \textit{HIV prevention: PrEP and PEP guidelines.} https://www.ashm.org.au/
    \item Healthdirect. \textit{HIV prevention and testing.} https://www.healthdirect.gov.au/hiv-aids
    \item Let's Talk About It and other community campaigns. \textit{PrEP and safe sex information.} https://www.health.nsw.gov.au/
    \item World Health Organization. \textit{HIV prevention.} https://www.who.int/
    \item UNAIDS. \textit{Prevention and treatment targets.} https://www.unaids.org/
    \item Centers for Disease Control and Prevention. \textit{HIV prevention basics.} https://www.cdc.gov/
\end{itemize}
